\section{Fixed even cycles have one logarithmic scale}\label{sec:cycles}

\begin{theorem}\label{thm:cycle}
Fix \(m\ge2\). On \(C_{2m}\) use the cap orders two and any seam
signing. With the normalization of \S\ref{sec:weights},
\[
 Z_r(M_{C_{2m},\epsilon})=
 \frac{\binom{2m-2}{m-1}}{\pi^{2m}}\, r^{2m}\log r
     +O_m(r^{2m}).
\]
There is no parity restriction on \(r\). In particular the leading
coefficient is positive and the leading logarithmic degree is one.
\end{theorem}

\begin{proof}
The adjacency signature is zero, and \(e=2m\) makes the parity
factor in \eqref{eq:finite} one. Label alternate defects \(a_i\)
and the intervening ones \(b_i\), with \(a_{m+1}=a_1\).
Put \(s_a=\min(a,r-a)\). The exact degree-two identity of
Paper~50 \cite{P50} is
\begin{equation}\label{eq:transfer}
 \sum_{\substack{1\le b<r\\ b\ {\rm odd}}}
 F_2(b/r)K(ab/r)K(cb/r)
 =\frac{r}{4ac}\min(s_a,s_c).
\end{equation}
An elementary proof also fixes the upper endpoint. The left
summand equals \(1/(4ac)\) times
\(\sin(2a\theta)\sin(2c\theta)/\sin^2\theta\), where
\(\theta=\pi b/(2r)\). On the odd grid \(a,c\) can be replaced
in the two sine numerators by \(s_a,s_c\), because
\(\sin(\pi b-2a\theta)=\sin(2a\theta)\).
Use \(\sin(2s\theta)/\sin\theta=2\sum_{j=1}^{s}
\cos((2j-1)\theta)\). On the full \(r\)-point midpoint grid
\(\theta=(b\pi)/(2r)\), \(b=1,3,\ldots,2r-1\), the cosine
modes of frequencies \(1,\ldots,r-1\) have norm squared \(r/2\)
and are mutually orthogonal by the finite geometric series.
Here \(2s_a-1,2s_c-1<r\). The full-grid sum of the sine product
ratios is \(2r\min(s_a,s_c)\); reflection about \(\pi/2\)
preserves the product, and the midpoint term, if present, is zero.
Take half and restore \(1/(4ac)\).

Each \(b_i\) can now be summed independently, leaving
\begin{align}
 Z_r(M_{C_{2m},\epsilon})
   &=\frac{2^m r^{2m}}{\pi^{2m}}T_m(r),\label{eq:Treln}\\
 T_m(r)
   &=\sum_{\substack{1\le a_i<r\\ a_i\ {\rm odd}}}
       \prod_{i=1}^m f_r(a_i)\min(s_{a_i},s_{a_{i+1}}),&
 f_r(a)&=\frac{\pi^2}{4r^2\sin^2(\pi a/(2r))}.\notag
\end{align}
Indeed \(I_r=(r/4)^m T_m\), because each \(a_i\) appears as
an endpoint of two transfers, so \(f_r(a)=a^{-2}F_2(a/r)\).
The reduced sum is positive termwise.

Let \(R=\lfloor r/2\rfloor\) and fold each label to \(s_a\), counting
the midpoint only once. Its weight is
\[
 w_r(s)=\sum_{\substack{1\le a<r,\ a\ {\rm odd}\\s_a=s}}f_r(a)
     =\frac{\Ind_{s\ {\rm odd}}}{s^2}+O(r^{-2}),\qquad
 0\le w_r(s)\le C s^{-2},\quad 1\le s\le R.
\]
The constants are uniform in parity: the lower preimage obeys the
first expansion by the bounded function
\(\csc^2 u-u^{-2}\) on \(0<u\le\pi/4\), while the upper preimage
has denominator \(4r^2\cos^2(\pi s/(2r))\) and contributes
\(O(r^{-2})\).

For \(s_{m+1}=s_1\), use the identity
\begin{equation}\label{eq:logkernel}
 \prod_i\frac{\min(s_i,s_{i+1})}{s_i^2}
  =\prod_i\frac1{s_i}
    \exp\!\left[-\frac12\sum_i|\log s_i-\log s_{i+1}|\right].
\end{equation}
For any fixed \(c>0\),
\(\sup_{a\ge1}\sum_{b\ge1}b^{-1}e^{-c|\log(b/a)|}<\infty\),
by splitting at \(a\). Holding one \(s_i\) fixed and deleting an
edge of the cycle, summation of the resulting path from its leaves
applies this bound repeatedly. In the telescoping replacement of the
\(m\) factors \(w_r(s_i)\), the erroneous site's cost relative to
\(s_i^{-2}\) is bounded by \(C r^{-2}s_i^2\), while the other
weights are bounded by constants times \(s_j^{-2}\). Therefore
each error term sums to \(O_m(r^{-2}\sum_{s\le R}s)=O_m(1)\).
This is an absolute error bound for the entire folded measure.

Compare the remaining sum on odd anchors with the integral of
the right side of \eqref{eq:logkernel} over cells
\(\prod_i[a_i,a_i+2]\). Over such a cell the logarithm of that
integrand changes by at most \(C_m\sum_i a_i^{-1}\).
Its ratio to the corner value is uniformly bounded. Use the same
path estimate for each inserted \(a_i^{-1}\) and
\(\sum_{a_i\ge1}a_i^{-2}<\infty\) to bound the sum-integral
error by \(O_m(1)\). With \(t_i=\log s_i\), this gives
\[
 T_m(r)=2^{-m}\int_{[0,L]^m}
          e^{-\frac12\sum_i|t_i-t_{i+1}|}\,\dd^m t+O_m(1),
 \qquad L=\log(R+O(1))=\log r+O(1).
\]
The factor \(2^{-m}\) is the exact leading density of the odd
lattice, not a fixed-parity assumption about the level.

Fix relative coordinates \(v_i=t_i-t_1\), \(v_1=v_{m+1}=0\).
Write \(D(v)=\max_i v_i-\min_i v_i\). For given \(v\), \(t_1\)
ranges on an interval of length \((L-D(v))_+\). Since the closed
variation \(\sum_i|v_i-v_{i+1}|\) is at least \(2D(v)\), the relative
weight has finite mass and finite first \(D\)-moment. Thus the
integral is \(L C_m+O_m(1)\), where
\[
 C_m=\int_{\RR^{m-1}} e^{-\frac12\sum_i|v_i-v_{i+1}|}\,\dd^{m-1}v .
\]
Equivalently \(C_m=k^{*m}(0)\), \(k(x)=e^{-|x|/2}\). In the
Fourier convention \(\widehat k(\xi)=\int e^{-i\xi x}k(x)\,\dd x\)
the transform is \((\xi^2+1/4)^{-1}\), so the beta integral gives
\[
 C_m=\frac1{2\pi}\int_{\RR}(\xi^2+1/4)^{-m}\,\dd\xi
 =\frac{2^{2m-2}\Gamma(m-\frac12)}{\sqrt\pi\,\Gamma(m)}
 =\binom{2m-2}{m-1}.
\]
Substituting \(T_m=2^{-m}C_m\log r+O_m(1)\) into
\eqref{eq:Treln} proves the assertion.
\end{proof}

\begin{corollary}[Common exponent and relative log colors]\label{cor:law}
Normalize the folded positive weight in \(T_m(r)\) to probability.
Then \((\log s_1/\log r,\,(\log(s_i/s_1))_{i=2}^m)\) converges
weakly to the product of the uniform law on \([0,1]\) and the density
\[
 C_m^{-1}\exp\!\left[-\tfrac12\sum_i|v_i-v_{i+1}|\right]\dd^{m-1}v.
\]
Thus all the retained alternating labels share the same limiting
exponent, whereas their relative logs retain a nontrivial law.
\end{corollary}
\begin{proof}
The weight replacement is \(O_m(1)\) in total variation. Away from
large relative differences, discard \(t_1\le M\) at cost \(O_m(M+1)\)
in the unnormalized weights and compare the remaining small cells
with the integral as before. With \(\max|v_i|\) bounded, their
diameters in the log variables tend uniformly to zero as
\(M\to\infty\), so this comparison also applies to bounded continuous
tests on fixed compact relative sets. The removed part vanishes
after division by \(\log r\) for fixed \(M\).
In the integral, set \(u=t_1/\log r\); its allowed interval tends to
\([0,1]\) uniformly for bounded \(v\).
Both discrete and continuous relative tails are uniformly
exponentially small after normalization: the factor
\(e^{-\sum|v_i-v_{i+1}|/4}\) bounds the cost of the other half of
the exponent whenever \(\max|v_i|\) is large, and the remaining
half is summed by the same fixed-anchor path bound. Let the
relative cutoff then tend to infinity.
\end{proof}

This law is a probability interpretation of the exact reduced
uninserted scalar. It does not assume that the folded labels in the
proof are jointly measurable Wilson angles. Even-cycle length is
visible here through the growth exponent \(2m\); no assertion compares
that exponent with every other type of backbone.
